\scanpage{63}
\mspage{106}
\begin{flushright}n° 371\end{flushright}

$F=\handwritten{\prod}_{n}E_n$. Montrons qu'on peut définir sur $F$ une structure d'anneau qui induise sur $E$ la structure d'anneau donnée ; en effet, tout élément de $F$ est de la forme $(u_n)$, où $u_n\handwritten{\in}E_n$ ; on définit sur $F$ une multiplication en posant $(u_n)(v_n)=(w_n)$, avec
\[
w_n=\handwritten{\sum}_{p+q=n}u_pv_q
\tag*{(28 bis)}
\]
ce qui a un sens quel que soit $n$, puisque le second membre n'a qu'un nombre fini de termes ; l'élément $1\handwritten{\in}E_{\unsure{o|0}}$ de $E$ est évidemment élément unité pour cette loi et l'associativité se vérifie immédiatement ; on a donc ainsi sur $E$ une structure d'anneau (commutatif si $E$ est commutatif). On peut de même définir sur $B=\handwritten{\prod}_{n}A_n$ une structure d'anneau commutatif ; enfin, on définit sur $F$ une structure de $B$-\uline{algèbre} de façon analogue : si $(\handwritten{\alpha}_n)$ est un élément de $B$, on posera $(\handwritten{\alpha}_n)(u_n)=(w'_n)$ avec
\[
w'_n=\handwritten{\sum}_{p+q=n}\handwritten{\alpha}_pu_q.
\tag*{(28 ter)}
\]

Nous dirons que $B$ (resp. $F$) est la \uline{complétion canonique} de l'anneau gradué $A$ (resp. de la $A$-algèbre graduée $E$).

\begin{quote}
\handwritten{*}Cette terminologie se justifie en considérant sur $B$ (resp. $E$) la topologie \uline{produit} des topologies discrètes sur les $A_n$ (resp. $E_n$) ; alors on voit aussitôt que la multiplication dans $F$ (resp. la loi externe de $B$-module sur $F$) est le \uline{prolongement par continuité} à $F\handwritten{\times}F$ (resp. $B\handwritten{\times}F$) de la multiplication $E\handwritten{\times}E\handwritten{\to}E$ (resp. de la loi externe $A\handwritten{\times}E\handwritten{\to}E$).\handwritten{*}
\end{quote}

Soit $E'$ une seconde $A$-algèbre graduée de type $Z$ à degrés positifs, et soit $F'$ sa complétion canonique. Tout homomorphisme $f:E\handwritten{\to}E'$ de $A$-algèbres graduées se prolonge alors de façon canonique en un homomorphisme $g:F\handwritten{\to}F'$ de $B$-algèbres ; il suffit de poser, pour tout $u=(u_n)\handwritten{\in}F$
\[
g(u)=(f(u_n))
\]
et on vérifie aussitôt qu'on définit bien ainsi un homomorphisme de $B$-algèbres.
% archive sheet number: 62
